\documentclass[a4paper,10pt]{article}
\usepackage[T1]{fontenc}
\usepackage[utf8]{inputenc}
% \usepackage{article}
\usepackage[margin=16mm]{geometry}
\usepackage{multicol}
\pagestyle{empty}
\setlength{\parindent}{0pt}
\setlength{\columnsep}{10mm}
\renewenvironment{itemize}{\begin{list}{\textbullet}{\setlength{\leftmargin}{1.3em}\setlength{\topsep}{3pt}\setlength{\itemsep}{2pt}\setlength{\parsep}{0pt}}}{\end{list}}
\newcommand{\tema}[1]{\par\addvspace{9pt}{\large\bfseries #1}\par\nobreak}

\begin{document}
\begin{center}
  {\LARGE\bfseries Sentrale tema for PPU 3210}\par
  \vspace{5pt}
  {\large Redegjør for sentrale observasjoner}
\end{center}

\small
Temaene og begrepene kan gi ulike «briller» for observasjon og hjelpe deg til å legge merke til sentrale elementer i videoen. Problemstillingen eller oppgaveteksten kan blant annet åpne for følgende temaer:

\begin{multicols}{2}
\raggedcolumns
\tema{Klasseledelse}
\begin{itemize}
  \item Strukturperspektivet
  \item Relasjonsperspektivet
  \item Fag- og læringsperspektivet
\end{itemize}

\tema{Organisering av undervisning}
\begin{itemize}
  \item Tilegnelse, utprøving, konsolidering (TUK)
  \item Faglig, emosjonell og organisatorisk støtte
\end{itemize}


\tema{Klasseromsdialoger}
\begin{itemize}
  \item Helklassedialog
  \item Gruppedialog
  \item Dialogisk undervisning
  \item Samtaleteknikker (\textit{talk moves})
  \item Samtaletyper (\textit{types of talk}):
    \begin{itemize}
      \item Akkumulerende
      \item Disputerende
      \item Utforskende/eksplorerende
    \end{itemize}
\end{itemize}

\tema{Læringsmiljø}
\begin{itemize}
  \item Relasjonsbygging
  \item Elevsentrert undervisning
  \item Sosial posisjonering
\end{itemize}

\columnbreak

\tema{Profesjonsfaglig digital kompetanse}
\begin{itemize}
  \item Primære og sekundære skoletekster
  \item Digitale kilder og KI
  \item Lærerrollen i digitaliserte klasserom
\end{itemize}

\tema{Noen viktige perspektiver på læring}
\begin{itemize}
  \item Læringsstrategier, selvregulering og dybdelæring
  \item Nærmeste utviklingssone
  \item Stillasbygging
  \item Modellering
  \item Elevenes forkunnskaper
  \item Begrepslæring
  \item Samarbeidslæring
\end{itemize}

\tema{Motivasjon}
\begin{itemize}
  \item Forventninger (mestring/relevans/kostnad)
  \item Behov for opplevd kompetanse, selvbestemmelse og tilhørighet
  \item Årsaksforklaring
  \item Målorientering
  \item Forsvar av selvbilde
  \item Indre motivasjon og autonom/kontrollert ytre motivasjon
\end{itemize}

\tema{Læringsteori}
\begin{itemize}
  \item Akkomodasjon og assimulasjon
      \begin{itemize}
          \item Tilpassning av kognitive skjema for å åpne "likevekt"
          \item Når assimulasjon ikke er mulig, så vil elevene streve etter akkomodasjon.
          \item Sammenligning med eksisterenede skjema og det vi erfarer
          \item Attribusjonsteori
          \item Hjelpløshet
      \end{itemize}
  \item Piaget
  \item Vigotsky
\end{itemize}

\tema{Matematikk}
\begin{itemize}
  \item Instrumentell og rasjonel forstålse
  \item Polyas Læringsstrategi
      \begin{itemize}
          \item Lage en plan. Gjennomføre. Kontrollere. Reflektere.
      \end{itemize}
  \item Akkumlasjon. Bruke det du har lært til nye problemstillinger (?)
  \item Sprik mellom eksisterende skjema og utfordringene som blir gitt. \implies  hjelpløshet
\end{itemize}
\end{multicols}

\vfill
{\footnotesize Kilde: «Oppsummering og videocaseeksamen, Høst 2026», slide 20: «Sentrale tema for PPU 3210». Oversikten er gjengitt med mindre språklige rettelser.}

\end{document}
